\begin{scanpage}{0197}{179}
length vibrator occurs in the same way as at the end of a semi-infinite conductor.

In the case of a passive vibrator ($z_1<z<z_2$), excited by the plane wave
\[
E_z^e=E_{0z}e^{iwz},\quad w=-k\cos\vartheta,\tag{28.13}
\]
the current also is represented in the form of the sum of waves (see [82])
\[
\begin{aligned}J(z)=S[&e^{iwz}-\psi_-(z-z_1)e^{iwz_1+ik(z-z_1)}\\
&-\psi_+(z_2-z)e^{iwz_2+ik(z_2-z)}+\overline A_1\psi(z-z_1)e^{ik(z-z_1)}\\
&+\overline A_2\psi(z_2-z)e^{ik(z_2-z)}],
\end{aligned}\tag{28.14}
\]
where the first term corresponds to the current excited by a plane wave in an infinitely long conductor. Its complex amplitude equals
\[
S=\frac{i\omega E_{0z}}{2k^2\sin^2\vartheta\ln\frac{2i}{\gamma ka\sin\vartheta}}.\tag{28.15}
\]
The second and third terms are primary edge waves arising as a consequence of the cut-off of the current $Se^{iwz}$. They are expressed in terms of the functions $\psi_+(z)$ and $\psi_-(z)$ which depend, in addition to the variable $z$ and the parameters $k$ and $a$, on the angle $\vartheta$. These functions satisfy the relationships
\[
\left.\begin{aligned}\psi_\pm(0)&=1,\quad\psi_\pm(\infty)=0,\\
\psi_+(z)\big|_{\vartheta=\pi}&=\psi_-(z)\big|_{\vartheta=0}=\psi(z),\\
\psi_+(z)\big|_{\vartheta=0}&=\psi_-(z)\big|_{\vartheta=\pi}=1.
\end{aligned}\right\}\tag{28.16}
\]
The initial values of the primary edge waves are such that their sum with the wave $Se^{iwz}$ gives a current equal to zero at the conductor ends.

The last two terms in Equation (28.14) correspond to secondary, tertiary, etc., edge waves, and have the same form as they do for a
\end{scanpage}
